\documentclass{article}

\usepackage[spanish]{babel}
\usepackage[utf8]{inputenc}
\usepackage{fancyhdr}
\usepackage{extramarks}
\usepackage{amsmath}
\usepackage{amsthm}
\usepackage{amsfonts}
\usepackage{tikz}
\usepackage{algorithm}
\usepackage{algpseudocode}
\usepackage{hyperref}
\usepackage{graphicx}
\usepackage{float}
\usepackage{natbib}

\usetikzlibrary{automata,positioning}

%
% Basic Document Settings
%

\topmargin=-0.45in
\evensidemargin=0in
\oddsidemargin=0in
\textwidth=6.5in
\textheight=9.0in
\headsep=0.25in

\linespread{1.1}

\pagestyle{fancy}
\lhead{\hmwkAuthorName}

\setlength{\headheight}{35pt}

\chead{
	\parbox{0.8\textwidth}{
		\centering
		\hmwkClass\ (\hmwkClassInstructor)\\
		\hmwkTitle
}}

\rhead{\firstxmark}
\lfoot{\lastxmark}
\cfoot{\thepage}

\renewcommand\headrulewidth{0.4pt}
\renewcommand\footrulewidth{0.4pt}

\setlength\parindent{0pt}

%
% Create Problem Sections
%

\newcommand{\enterProblemHeader}[2]{
	\nobreak\extramarks{}{Problema {#1} continúa en la sig pág\ldots}\nobreak{}
	\nobreak\extramarks{Problema \arabic{#1} (continúa)}{Problema \arabic{#1} continúa en la sig pág\ldots}\nobreak{}
}

\newcommand{\exitProblemHeader}[2]{
	\nobreak\extramarks{Problema \arabic{#1}}{Problema \arabic{#1} continúa en la sig pág\ldots}\nobreak{}
	\stepcounter{#1}
	\nobreak\extramarks{Problema \arabic{#1}}{}\nobreak{}
}

\setcounter{secnumdepth}{0}
\newcounter{partCounter}
\newcounter{homeworkProblemCounter}
\setcounter{homeworkProblemCounter}{1}
\nobreak\extramarks{Problema \arabic{homeworkProblemCounter}}{}\nobreak{}

%
% Homework Problem Environment
%
% This environment takes an optional argument. When given, it will adjust the
% problem counter. This is useful for when the problems given for your
% assignment aren't sequential. See the last 3 problems of this template for an
% example.
%
\newenvironment{homeworkProblem}[2][-1]{
	\ifnum#1>0
	\setcounter{homeworkProblemCounter}{#1}
	\fi
	\def\currentProblemTitle{#2}
	\section{Problema \arabic{homeworkProblemCounter}: #2}
	\setcounter{partCounter}{1}
	\enterProblemHeader{homeworkProblemCounter}{#2}
}{
	\exitProblemHeader{homeworkProblemCounter}{\currentProblemTitle}
}
%
% Homework Details
%   - Title
%   - Due date
%   - Class
%   - Section/Time
%   - Instructor
%   - Author
%

\newcommand{\hmwkTitle}{Completitud del Tangent Bug\ \#1}
\newcommand{\hmwkDueDate}{21 de mayo de 2026}
\newcommand{\hmwkClass}{Planificación de movimientos de robots}
\newcommand{\hmwkClassInstructor}{Rafael Murrieta Cid}
\newcommand{\hmwkAuthorName}{\textbf{Andrés Alemán}}

%
% Title Page
%

\title{
	\vspace{2in}
	\textmd{\textbf{\hmwkClass:\ \hmwkTitle}}\\
	\normalsize\vspace{0.1in}\small{Due\ on\ \hmwkDueDate}\\
	\vspace{0.1in}\large{\textit{\hmwkClassInstructor}}
	\vspace{3in}
}

\author{\hmwkAuthorName}
\date{}

\renewcommand{\part}[1]{\textbf{\large Parte \Alph{partCounter}}\stepcounter{partCounter}\\}

%
% Various Helper Commands
%

% Useful for algorithms
\newcommand{\alg}[1]{\textsc{\bfseries \footnotesize #1}}

% For derivatives
\newcommand{\deriv}[1]{\frac{\mathrm{d}}{\mathrm{d}x} (#1)}

% For partial derivatives
\newcommand{\pderiv}[2]{\frac{\partial}{\partial #1} (#2)}

% Integral dx
\newcommand{\dx}{\mathrm{d}x}

% Alias for the Solution section header
\newcommand{\solution}{\textbf{\large Solución}}

% Probability commands: Expectation, Variance, Covariance, Bias
\newcommand{\E}{\mathrm{E}}
\newcommand{\Var}{\mathrm{Var}}
\newcommand{\Cov}{\mathrm{Cov}}
\newcommand{\Bias}{\mathrm{Bias}}

\begin{document}
	
%	\maketitle
	
%	\pagebreak
	
	\begin{homeworkProblem}{Determinar si Tangent Bug es completo o no}
		
		Antes de determinar la completitud, se define el algoritmo de acuerdo a \citet{choset2005principles}, quienes describen Tangent Bug como:
					
			\begin{algorithm}
				\caption{Tangent Bug Algorithm}
				\textbf{Input:} A point robot with a range sensor \\
				\textbf{Output:} A path to $q_{goal}$ or a conclusion no such path exists
				\hrule
				
				\begin{algorithmic}[1]
					\While{True}
					\Repeat
					\State Continuously move toward the point $n \in \{T, O_i\}$ which minimizes 
					$d(x,n) + d(n,q_{goal})$
					\Until{the goal is encountered \textbf{or} the direction that minimizes $d(x,n) + d(n,q_{goal})$ 
						begins to increase $d(x,q_{goal})$ (i.e., a local minimum of $d(\cdot, q_{goal})$ is detected)}
					
					\State Choose a boundary-following direction which continues in the same direction as the most recent motion-to-goal direction
					
					\Repeat
					\State Continuously update $d_{reach}$, $d_{followed}$, and $\{O_i\}$
					\State Continuously move toward $n \in \{O_i\}$ that is in the chosen boundary direction
					\Until{
						the goal is reached \textbf{or} 
						the robot completes a cycle around the obstacle (goal not achievable) \textbf{or} 
						$d_{reach} < d_{followed}$
					}
					\EndWhile
				\end{algorithmic}
			\end{algorithm}
		
		En donde:
		
		\begin{itemize}
			\item $x \in \mathbb{R}^2$: posición actual del robot en el espacio de trabajo.
			
			\item $q_{goal} \in \mathbb{R}^2$: configuración objetivo que el robot debe alcanzar.
			
			\item $d(a, b) = \|a - b\|$: distancia euclidiana entre dos puntos $a$ y $b$ en el espacio de trabajo.
			
			\item $O_i$: punto de discontinuidad en el rango del sensor asociado al $i$-ésimo obstáculo, correspondiente a los extremos (esquinas) detectados por el robot.
			
			\item $T$: conjunto de puntos tangentes visibles desde la posición actual $x$, definidos como los puntos $O_i$ desde los cuales la línea de visión al robot es tangente al obstáculo correspondiente.
			
			\item $n \in \{T \cup \{O_i\}\}$: punto candidato hacia el cual el robot evalúa moverse; puede ser un punto tangente o un punto de discontinuidad visible.
			
			\item $d(x, n) + d(n, q_{goal})$: función de costo que estima la distancia total del camino pasando por el punto candidato $n$; el robot se mueve hacia el $n$ que minimiza esta expresión.
			
			\item $d_{followed}$: distancia euclidiana $d(x, q_{goal})$ registrada en el momento en que el robot inicia la fase de seguimiento de frontera (\textit{boundary-following}). Sirve como referencia para evaluar el progreso.
			
			\item $d_{reach}$: menor valor de $d(n, q_{goal})$ observado sobre todos los puntos tangentes $n \in T$ durante la fase de seguimiento de frontera. Representa la distancia mínima a la meta que el robot podría alcanzar desde su trayectoria actual.
		\end{itemize}
		
		Se dice que \textit{Tangent Bug} es completo si, para cualquier configuración 
		inicial $x$ y objetivo $q_{goal}$, el algoritmo termina en tiempo finito 
		garantizando uno de los siguientes casos:
		
		\begin{itemize}
			\item El robot alcanza $q_{goal}$.
			\item El robot determina que $q_{goal}$ es inalcanzable.
		\end{itemize}
		
		Esto se garantiza por las condiciones de terminación de la fase de 
		\textit{boundary-following}: si $d_{reach} < d_{followed}$, el robot 
		progresa hacia la meta; de lo contrario, al completar un ciclo alrededor 
		del obstáculo sin mejora, el algoritmo concluye que no existe camino libre.
		
	\end{homeworkProblem}
	
	\bibliographystyle{plainnat}
	\bibliography{biblio}
\end{document}
